\subsection{Fredholm mappings and Riesz mappings between Fréchet spaces}

PROPOSITION 11. --- Let E and F be Fréchet spaces. For a continuous linear map u from E into F to be a Fredholm map, it is necessary and sufficient that its kernel and cokernel be of finite dimension.

This follows from the characterization of Fredholm maps given by condition (ii) of Prop. 2 in III, p. 42, and from the following lemma:

Lemma 6. --- Let E and F be Fréchet spaces and \(u \in \mathcal{L}(\mathrm{E};\mathrm{F})\) . We assume that the image of \(u\) has finite codimension in F. Then the image of \(u\) is closed in F and \(u\) is a strict morphism.

Let G be a vector subspace of F supplementary to \(\operatorname{Im}(u)\) ; the subspace G is closed (EVT, I, p. 14, cor. 1), hence the quotient space F/G is a Fréchet space (TG, IX, p. 25, prop. 4). Let \(\pi: F \to F/G\) be the canonical surjection. The map \(\pi \circ u\) is surjective, hence strict (EVT, I, p. 19, cor. 3). The assertion then follows from cor. 2 of III, p. 40.

PROPOSITION 12. --- Let E and F be Fréchet spaces. Endow their duals E' and F' with the weak topology, the topology of compact convergence, or the topology of bounded convergence. Let \(u \in \mathcal{L}(\mathrm{E};\mathrm{F})\) . For u to be a Fredholm map from E to F, it is necessary and sufficient that \({}^{t}u\) be a Fredholm map from F' to E'.

The condition is necessary (III, p. 43, no. 3). Let us prove that it is sufficient. If \(^{t}u\) is a Fredholm map, it is a strict morphism (III, p. 42, prop. 2), hence u is a strict morphism with closed image (EVT, IV, p. 28, th. 1 and p. 29, cor. 3). Taking into account the canonical isomorphisms Coker( \(^{t}u\) ) → Ker(u)' and Ker( \(^{t}u\) ) → Coker(u)' (EVT, IV, p. 27, prop. 2), the vector spaces Ker(u) and Coker(u) are finite-dimensional and u is a Fredholm operator (prop. 2 of III, p. 42).

PROPOSITION 13. --- Let E be a Fréchet space and \(u \in \mathcal{L}(\mathrm{E})\) .

\begin{enumerate}
\def\labelenumi{\alph{enumi})}
\item
  The endomorphism u is a Riesz endomorphism if and only if its nilespace N is finite-dimensional and its conilespace I is of finite codimension;
\item
 Endow E' with the topology of bounded convergence, the topology of compact convergence, or the weak topology. If \(^{t}\) u is a Riesz endomorphism of E', then u is a Riesz endomorphism of E.
\end{enumerate}
a) Suppose that the nilespace N is finite-dimensional and that the conilespace I is of finite codimension. Since \(\operatorname{Ker}(u) \subset \operatorname{N}\) and \(\operatorname{I} \subset \operatorname{Im}(u)\) , the kernel of u is finite-dimensional and the image of u is of finite codimension. By the definition and proposition 11, the map u is a Riesz endomorphism. The converse follows from the definition.

b) By proposition 12, the hypothesis implies that \(u\) is a Fredholm endomorphism of index \(\text{ind}(^t u) = -\text{ind}(u) = 0\) ( \(n^o 3\) , formula (4)). Let \(n \in \mathbf{N}\) . Since the image of \(u^n\) is closed and has as its orthogonal the kernel of \(^t u^n\) (EVT, IV, p. 27, prop. 2), the sequence \((\text{Im}(u^n))\) is stationary, and \(u\) is therefore a Riesz endomorphism.


